\section{JASPでベイズ分析を体験する}
\label{b28_bayesian_with_jasp}

\subsection{授業内容}

\subsubsection{概要}
前コマまでにベイズ的な判断の方法を学んだが，これを実際に行うにはソフトウェアの助けが要る。プログラミングの習得を前提とせずにベイズ的アプローチを試せる道具として，本コマではJASPを用いる。JASPはオープンソースの統計ソフトウェアで，心理学研究に特化した機能と直感的なユーザーインターフェイスを提供し，頻度主義的方法とベイズ的方法を同じメニューから選べる点に特徴がある。従来通りの分析をしつつ，ベイズでやると何がどう変わるのかを比較しながら進められるわけである。本コマでは，JASPの基本的な操作方法から始め，記述統計と可視化，回帰分析におけるベイズ推定，そしてベイズファクターを用いた検定について実践的に学ぶ。これにより，学生はベイズ統計学の理論と実践の両面から理解を深め，心理学研究におけるデータ分析能力を向上させることが期待される。

\subsubsection{コマ主題細目}
\begin{description}
	\item[JASPの導入とインターフェイス] JASPはオープンソースの統計ソフトウェアで，心理学研究のための統計分析機能と直感的なユーザーインターフェイスを提供する。まず，JASPのウェブサイト（https://jasp-stats.org/）からソフトウェアをダウンロードし，各自のPCにインストールする方法を学ぶ。続いて，JASPの基本的なインターフェイスと機能（メニュー構成，データビュー，分析ビュー，結果ビュー）について理解する。JASPの特徴として，オープンサイエンスをサポートする機能（分析結果の再現性，OSFとの連携）や，従来の頻度主義的分析とベイズ的分析を同じメニューから実行できる利便性について学ぶ。また，サンプルデータの読み込み方，外部データ（CSV，SPSSなど）のインポート方法，変数の種類の設定方法についても実習を通じて習得する。
	      \begin{flushright}
		      $\to$ \textcite{YUnaShimizu2022}のP.1--15.
	      \end{flushright}

	\item[記述統計と可視化] JASPを用いた基本的なデータ分析として，記述統計と可視化の方法を学ぶ。「Descriptives」メニューを使って，各変数の基本統計量（平均，中央値，標準偏差，四分位数など）を計算し，結果の読み方を理解する。また，ヒストグラム，箱ひげ図などのグラフを作成し，データの分布特性を視覚的に把握する方法を学ぶ。これらの基本操作を通じて，JASPのワークフローに慣れ，次のステップであるより高度な分析への準備を整える。
		\begin{flushright}
		      $\to$ \textcite{YUnaShimizu2022}のP.1--15.
	      \end{flushright}

	\item[JASPによる相関・回帰分析] JASPを用いた回帰分析を，頻度主義的アプローチとベイズ的アプローチの両方で実施する。「回帰」メニューの「線形回帰」を伝統的手法とベイズ的手法，それぞれ使用して分析を行い，結果の違いを比較する。ベイズ回帰分析では，回帰係数の事後分布，確信区間，モデル比較のためのベイズファクターが出力される。

	\item[ベイズファクターを使った検定] JASPを使った平均値差の検定は，ベイズファクターを用いたモデル比較の形で行われる。例えば，2群の平均値差を検定する場合，帰無仮説モデル（平均値が等しい）と対立仮説モデル（平均値が異なる）を設定し，ベイズファクター$BF_{10}$を計算する。$BF_{10}$の値に基づいて，どちらのモデルがデータによりよく適合しているかを判断する。JASPでは，この分析を簡単に実行でき，結果の解釈方法も学ぶ。

\end{description}
\subsubsection{キーワード}
\begin{itemize}
	\item JASP
	\item GUIによるベイズ分析
	\item ベイズ回帰分析
	\item ベイズファクターによる検定
\end{itemize}
\subsection{授業情報}
\paragraph{コマの展開方法}
講義
\paragraph{標準シラバスにおける位置づけ}
科目番号5；心理学統計法；(1)心理学で用いられる統計手法；J より高度な記述や推定を目指して
\subsubsection{予習・復習課題}
\paragraph{予習}
前回までに学んだベイズ的判断の方法（確信区間，ROPE，ベイズファクター）について復習しておくこと。特に以下の点を確認しておくとよい：
\begin{itemize}
  \item ベイズファクター$BF_{10}$の定義と解釈基準（1--3: わずかな証拠，3--10: 中程度，10以上: 強力な証拠）
  \item 確信区間とROPEに基づく「実質的等価性」の判断方法
  \item 頻度主義的な回帰分析の基本（回帰係数，決定係数$R^2$，$t$検定など）
\end{itemize}
また，JASPを各自のPCにインストールしておくこと（https://jasp-stats.org/）。インストールが完了したら，起動して基本的なインターフェイスを確認しておくと，授業での実習がスムーズに進む。参考文献として\textcite{YUnaShimizu2022}のP.1--15を事前に読んでおくことが望ましい。

\paragraph{復習}
授業で学んだ内容を定着させるために，以下の点について理解を深めておくこと：
\begin{enumerate}
  \item JASPの基本操作（データの読み込み，変数の設定，分析メニューの選択，結果の読み方）を実際に操作して確認する
  \item 回帰モデルのベイズ表現について，パラメータ$\beta_0$，$\beta_1$，$\sigma^2$に事前分布を設定する意味を説明できるようにする
  \item 頻度主義的回帰分析とベイズ回帰分析の結果の違い（点推定値，区間推定，仮説検証の方法）を比較し，それぞれの特徴を理解する
  \item JASPで出力されるベイズファクターの読み方と解釈方法を，具体的な数値例を用いて説明できるようにする
  \item サンプルデータを用いて，JASPでの記述統計，可視化，回帰分析を一通り実行し，操作に習熟しておく
\end{enumerate}
次回はGUIでは手の届かない領域に進む。JASPのメニューにない分析をやりたくなったときにどうするか，という問題意識を持っておくと，次コマの内容が理解しやすい。
